% Seccion del TFM: procedencia de cada kernel. Incluir con % Seccion del TFM: procedencia de cada kernel. Incluir con % Seccion del TFM: procedencia de cada kernel. Incluir con % Seccion del TFM: procedencia de cada kernel. Incluir con \input{docs/TFM/procedencia_kernels.tex}.
% Requiere: \usepackage{booktabs,tabularx}. Version Markdown: docs/TFM/procedencia_kernels.md.

\section{Procedencia de los kernels}
\label{sec:procedencia}

Esta sección indica el origen de cada kernel: si es código oficial sin modificar, una
adaptación de un ejemplo público o una implementación propia sobre una API documentada.
La procedencia se ha verificado comparando el código con la fuente citada, en la versión
indicada en la tabla~\ref{tab:fuentes}. Se distinguen tres tipos de elaboración:
\textbf{oficial sin modificar} (se llama al algoritmo de la fuente tal cual, con la
envoltura mínima para medirlo), \textbf{adaptado} (parte de un ejemplo público, con los
cambios que se detallan) y \textbf{propio} (diseño propio sobre la API documentada).

\subsection{Kernels compute-bound}

\begin{table}[htbp]
  \centering\small
  \begin{tabularx}{\linewidth}{llX}
    \toprule
    kernel (\texttt{src/CBKernels/triton/}) & elaboración & fuente \\
    \midrule
    \texttt{matmul.py} (baseline) & adaptado & tutorial de Triton \texttt{03-matrix-multiplication}~[1] \\
    \texttt{matmul\_blockptr.py} & adaptado & baseline + \texttt{tl.make\_block\_ptr}~[1] \\
    \texttt{matmul\_tma.py} & adaptado & tutorial de Triton \texttt{09-persistent-matmul} (variantes TMA)~[1] \\
    \texttt{matmul\_fp8.py} & adaptado & baseline con entradas FP8 e4m3~[1] \\
    \texttt{matmul\_fp8\_tma.py} & propio & combinación de \texttt{matmul\_fp8} y \texttt{matmul\_tma} \\
    \midrule
    \texttt{helion/matmul.py}$^\dagger$ & adaptado & ejemplo oficial de Helion \texttt{examples/matmul.py}~[3] \\
    \texttt{cutlass/matmul\_cute.py}$^\dagger$ & oficial sin modificar & CuTe DSL \texttt{blackwell\_geforce/kernel/dense\_gemm/dense\_gemm.py}~[2] \\
    \bottomrule
  \end{tabularx}
  \caption{Procedencia de los kernels de matmul. $^\dagger$Ruta relativa a \texttt{src/CBKernels/}.}
  \label{tab:procedencia-cb}
\end{table}

El baseline \texttt{matmul.py} comparte 27 de sus 36 líneas distintas con el
\texttt{matmul\_kernel} del tutorial~03. Coinciden el orden de programas agrupado para la
L2 (\texttt{GROUP\_SIZE\_M}), la aritmética de punteros y el bucle en K con
\texttt{tl.dot} y acumulación en fp32. Son propios el espacio de autotuning y la
comprobación en el PTX de que se usan tensor cores. \texttt{matmul\_tma.py} toma del
tutorial~09 los descriptores de tensor creados dentro del kernel
(\texttt{tl.make\_tensor\_descriptor}) y el \emph{allocator} obligatorio
(\texttt{triton.set\_allocator}). A diferencia del tutorial, no es persistente: conserva
la rejilla del baseline para aislar el efecto de TMA. \texttt{matmul\_blockptr.py} se
mantiene como intento documentado de activar TMA con \emph{block pointers}, que en
sm\_121 bajan a \texttt{cp.async} clásico (sección~\ref{sec:tma}). El matmul de Helion
es el del ejemplo oficial sin el \emph{epilogue}. El de CUTLASS ejecuta sin modificar el
ejemplo de CuTe DSL para Blackwell GeForce (\texttt{Sm120GemmKernel}: fp16, TMA y pipeline de
varias etapas). Solo es propia la envoltura que le pasa las matrices de PyTorch sin copiarlas.
No admite bf16.

\subsection{Kernels memory-bound: RMSNorm}

\begin{table}[htbp]
  \centering\small
  \begin{tabularx}{\linewidth}{lllX}
    \toprule
    kernel (\texttt{src/MBKernels/}) & DSL & elaboración & fuente \\
    \midrule
    \texttt{triton/rmsnorm\_baseline.py} & Triton & adaptado & tutorial de Triton \texttt{05-layer-norm}~[1] \\
    \texttt{triton\_tlx/rmsnorm.py} & Triton + TLX & propio & API de µTLX (\texttt{triton-ext})~[4] \\
    \texttt{gluon/rmsnorm.py} & Gluon & propio & API de Gluon, tutorial \texttt{gluon/02-layouts}~[1] \\
    \texttt{helion/rmsnorm.py} & Helion & adaptado & ejemplo oficial \texttt{examples/rms\_norm.py}~[3] \\
    \texttt{cutlass/rmsnorm\_ext.cu} & CUTLASS & oficial sin modificar & \texttt{cutlass::rmsnorm} (\texttt{device\_rmsnorm.h})~[2] \\
    \bottomrule
  \end{tabularx}
  \caption{Procedencia de las implementaciones de RMSNorm.}
  \label{tab:procedencia-mb}
\end{table}

El RMSNorm en Triton adapta el LayerNorm del tutorial~05: un programa por fila,
\texttt{MAX\_FUSED\_SIZE}~$=65536$ y la heurística de warps \texttt{BLOCK\_SIZE // 256}
(hasta 16 warps en lugar de 8). Se elimina la resta de la media. La versión Gluon repite
el algoritmo, pero fija el reparto de la fila entre hilos con un \texttt{BlockedLayout}
explícito. La de Helion es el \texttt{rms\_norm\_fwd} oficial sin la salida
\texttt{inv\_rms}, que solo necesita el \emph{backward}. La de CUTLASS invoca la función
oficial \texttt{cutlass::rmsnorm} sin modificarla: en fp16 usa el kernel vectorizado con
\texttt{float4} y en bf16 el escalar. Solo la envoltura como extensión de PyTorch es
propia. La versión TLX es un diseño propio (kernel persistente con \emph{prefetch} de
filas a memoria compartida mediante \texttt{cp.async}) que sigue el patrón de los tests y
tutoriales de triton-ext. Con \texttt{triton-utlx}~3.8.0.post1 no compila por un error del
plugin: crea \texttt{ttg.async\_copy\_global\_to\_local} sin \texttt{operandSegmentSizes},
y el fallo se reproduce también con el ejemplo mínimo del propio repositorio.

Como referencia de rendimiento y de corrección se usan implementaciones de PyTorch, no
kernels propios: \texttt{torch.matmul} (cuBLAS), \texttt{torch.\_scaled\_mm} (cuBLASLt,
FP8), \texttt{to\_sparse\_semi\_structured} (cuSPARSELt 2:4, con poda 2:4 propia) y
\texttt{torch.nn.functional.rms\_norm}.

\begin{table}[htbp]
  \centering\small
  \begin{tabularx}{\linewidth}{lXll}
    \toprule
    & proyecto & versión & licencia \\
    \midrule
    {[1]} & Triton (\texttt{triton-lang/triton}, \texttt{python/tutorials/}) & v3.8.0 (\texttt{c01b677}) & MIT \\
    {[2]} & CUTLASS (\texttt{NVIDIA/cutlass}) & v4.8.0 (\texttt{098de2a}) & BSD-3-Clause \\
    {[3]} & Helion (\texttt{pytorch/helion}) & \texttt{main} (\texttt{0f3b242}) & BSD-3-Clause \\
    {[4]} & µTLX (\texttt{triton-lang/triton-ext}, \texttt{triton-utlx} 3.8.0.post1) & \texttt{main} (\texttt{3c2bb3c}) & MIT \\
    \bottomrule
  \end{tabularx}
  \caption{Fuentes de los kernels. Todas las licencias permiten reutilizar y modificar el código conservando el aviso de copyright.}
  \label{tab:fuentes}
\end{table}
.
% Requiere: \usepackage{booktabs,tabularx}. Version Markdown: docs/TFM/procedencia_kernels.md.

\section{Procedencia de los kernels}
\label{sec:procedencia}

Esta sección indica el origen de cada kernel: si es código oficial sin modificar, una
adaptación de un ejemplo público o una implementación propia sobre una API documentada.
La procedencia se ha verificado comparando el código con la fuente citada, en la versión
indicada en la tabla~\ref{tab:fuentes}. Se distinguen tres tipos de elaboración:
\textbf{oficial sin modificar} (se llama al algoritmo de la fuente tal cual, con la
envoltura mínima para medirlo), \textbf{adaptado} (parte de un ejemplo público, con los
cambios que se detallan) y \textbf{propio} (diseño propio sobre la API documentada).

\subsection{Kernels compute-bound}

\begin{table}[htbp]
  \centering\small
  \begin{tabularx}{\linewidth}{llX}
    \toprule
    kernel (\texttt{src/CBKernels/triton/}) & elaboración & fuente \\
    \midrule
    \texttt{matmul.py} (baseline) & adaptado & tutorial de Triton \texttt{03-matrix-multiplication}~[1] \\
    \texttt{matmul\_blockptr.py} & adaptado & baseline + \texttt{tl.make\_block\_ptr}~[1] \\
    \texttt{matmul\_tma.py} & adaptado & tutorial de Triton \texttt{09-persistent-matmul} (variantes TMA)~[1] \\
    \texttt{matmul\_fp8.py} & adaptado & baseline con entradas FP8 e4m3~[1] \\
    \texttt{matmul\_fp8\_tma.py} & propio & combinación de \texttt{matmul\_fp8} y \texttt{matmul\_tma} \\
    \midrule
    \texttt{helion/matmul.py}$^\dagger$ & adaptado & ejemplo oficial de Helion \texttt{examples/matmul.py}~[3] \\
    \texttt{cutlass/matmul\_cute.py}$^\dagger$ & oficial sin modificar & CuTe DSL \texttt{blackwell\_geforce/kernel/dense\_gemm/dense\_gemm.py}~[2] \\
    \bottomrule
  \end{tabularx}
  \caption{Procedencia de los kernels de matmul. $^\dagger$Ruta relativa a \texttt{src/CBKernels/}.}
  \label{tab:procedencia-cb}
\end{table}

El baseline \texttt{matmul.py} comparte 27 de sus 36 líneas distintas con el
\texttt{matmul\_kernel} del tutorial~03. Coinciden el orden de programas agrupado para la
L2 (\texttt{GROUP\_SIZE\_M}), la aritmética de punteros y el bucle en K con
\texttt{tl.dot} y acumulación en fp32. Son propios el espacio de autotuning y la
comprobación en el PTX de que se usan tensor cores. \texttt{matmul\_tma.py} toma del
tutorial~09 los descriptores de tensor creados dentro del kernel
(\texttt{tl.make\_tensor\_descriptor}) y el \emph{allocator} obligatorio
(\texttt{triton.set\_allocator}). A diferencia del tutorial, no es persistente: conserva
la rejilla del baseline para aislar el efecto de TMA. \texttt{matmul\_blockptr.py} se
mantiene como intento documentado de activar TMA con \emph{block pointers}, que en
sm\_121 bajan a \texttt{cp.async} clásico (sección~\ref{sec:tma}). El matmul de Helion
es el del ejemplo oficial sin el \emph{epilogue}. El de CUTLASS ejecuta sin modificar el
ejemplo de CuTe DSL para Blackwell GeForce (\texttt{Sm120GemmKernel}: fp16, TMA y pipeline de
varias etapas). Solo es propia la envoltura que le pasa las matrices de PyTorch sin copiarlas.
No admite bf16.

\subsection{Kernels memory-bound: RMSNorm}

\begin{table}[htbp]
  \centering\small
  \begin{tabularx}{\linewidth}{lllX}
    \toprule
    kernel (\texttt{src/MBKernels/}) & DSL & elaboración & fuente \\
    \midrule
    \texttt{triton/rmsnorm\_baseline.py} & Triton & adaptado & tutorial de Triton \texttt{05-layer-norm}~[1] \\
    \texttt{triton\_tlx/rmsnorm.py} & Triton + TLX & propio & API de µTLX (\texttt{triton-ext})~[4] \\
    \texttt{gluon/rmsnorm.py} & Gluon & propio & API de Gluon, tutorial \texttt{gluon/02-layouts}~[1] \\
    \texttt{helion/rmsnorm.py} & Helion & adaptado & ejemplo oficial \texttt{examples/rms\_norm.py}~[3] \\
    \texttt{cutlass/rmsnorm\_ext.cu} & CUTLASS & oficial sin modificar & \texttt{cutlass::rmsnorm} (\texttt{device\_rmsnorm.h})~[2] \\
    \bottomrule
  \end{tabularx}
  \caption{Procedencia de las implementaciones de RMSNorm.}
  \label{tab:procedencia-mb}
\end{table}

El RMSNorm en Triton adapta el LayerNorm del tutorial~05: un programa por fila,
\texttt{MAX\_FUSED\_SIZE}~$=65536$ y la heurística de warps \texttt{BLOCK\_SIZE // 256}
(hasta 16 warps en lugar de 8). Se elimina la resta de la media. La versión Gluon repite
el algoritmo, pero fija el reparto de la fila entre hilos con un \texttt{BlockedLayout}
explícito. La de Helion es el \texttt{rms\_norm\_fwd} oficial sin la salida
\texttt{inv\_rms}, que solo necesita el \emph{backward}. La de CUTLASS invoca la función
oficial \texttt{cutlass::rmsnorm} sin modificarla: en fp16 usa el kernel vectorizado con
\texttt{float4} y en bf16 el escalar. Solo la envoltura como extensión de PyTorch es
propia. La versión TLX es un diseño propio (kernel persistente con \emph{prefetch} de
filas a memoria compartida mediante \texttt{cp.async}) que sigue el patrón de los tests y
tutoriales de triton-ext. Con \texttt{triton-utlx}~3.8.0.post1 no compila por un error del
plugin: crea \texttt{ttg.async\_copy\_global\_to\_local} sin \texttt{operandSegmentSizes},
y el fallo se reproduce también con el ejemplo mínimo del propio repositorio.

Como referencia de rendimiento y de corrección se usan implementaciones de PyTorch, no
kernels propios: \texttt{torch.matmul} (cuBLAS), \texttt{torch.\_scaled\_mm} (cuBLASLt,
FP8), \texttt{to\_sparse\_semi\_structured} (cuSPARSELt 2:4, con poda 2:4 propia) y
\texttt{torch.nn.functional.rms\_norm}.

\begin{table}[htbp]
  \centering\small
  \begin{tabularx}{\linewidth}{lXll}
    \toprule
    & proyecto & versión & licencia \\
    \midrule
    {[1]} & Triton (\texttt{triton-lang/triton}, \texttt{python/tutorials/}) & v3.8.0 (\texttt{c01b677}) & MIT \\
    {[2]} & CUTLASS (\texttt{NVIDIA/cutlass}) & v4.8.0 (\texttt{098de2a}) & BSD-3-Clause \\
    {[3]} & Helion (\texttt{pytorch/helion}) & \texttt{main} (\texttt{0f3b242}) & BSD-3-Clause \\
    {[4]} & µTLX (\texttt{triton-lang/triton-ext}, \texttt{triton-utlx} 3.8.0.post1) & \texttt{main} (\texttt{3c2bb3c}) & MIT \\
    \bottomrule
  \end{tabularx}
  \caption{Fuentes de los kernels. Todas las licencias permiten reutilizar y modificar el código conservando el aviso de copyright.}
  \label{tab:fuentes}
\end{table}
.
% Requiere: \usepackage{booktabs,tabularx}. Version Markdown: docs/TFM/procedencia_kernels.md.

\section{Procedencia de los kernels}
\label{sec:procedencia}

Esta sección indica el origen de cada kernel: si es código oficial sin modificar, una
adaptación de un ejemplo público o una implementación propia sobre una API documentada.
La procedencia se ha verificado comparando el código con la fuente citada, en la versión
indicada en la tabla~\ref{tab:fuentes}. Se distinguen tres tipos de elaboración:
\textbf{oficial sin modificar} (se llama al algoritmo de la fuente tal cual, con la
envoltura mínima para medirlo), \textbf{adaptado} (parte de un ejemplo público, con los
cambios que se detallan) y \textbf{propio} (diseño propio sobre la API documentada).

\subsection{Kernels compute-bound}

\begin{table}[htbp]
  \centering\small
  \begin{tabularx}{\linewidth}{llX}
    \toprule
    kernel (\texttt{src/CBKernels/triton/}) & elaboración & fuente \\
    \midrule
    \texttt{matmul.py} (baseline) & adaptado & tutorial de Triton \texttt{03-matrix-multiplication}~[1] \\
    \texttt{matmul\_blockptr.py} & adaptado & baseline + \texttt{tl.make\_block\_ptr}~[1] \\
    \texttt{matmul\_tma.py} & adaptado & tutorial de Triton \texttt{09-persistent-matmul} (variantes TMA)~[1] \\
    \texttt{matmul\_fp8.py} & adaptado & baseline con entradas FP8 e4m3~[1] \\
    \texttt{matmul\_fp8\_tma.py} & propio & combinación de \texttt{matmul\_fp8} y \texttt{matmul\_tma} \\
    \midrule
    \texttt{helion/matmul.py}$^\dagger$ & adaptado & ejemplo oficial de Helion \texttt{examples/matmul.py}~[3] \\
    \texttt{cutlass/matmul\_cute.py}$^\dagger$ & oficial sin modificar & CuTe DSL \texttt{blackwell\_geforce/kernel/dense\_gemm/dense\_gemm.py}~[2] \\
    \bottomrule
  \end{tabularx}
  \caption{Procedencia de los kernels de matmul. $^\dagger$Ruta relativa a \texttt{src/CBKernels/}.}
  \label{tab:procedencia-cb}
\end{table}

El baseline \texttt{matmul.py} comparte 27 de sus 36 líneas distintas con el
\texttt{matmul\_kernel} del tutorial~03. Coinciden el orden de programas agrupado para la
L2 (\texttt{GROUP\_SIZE\_M}), la aritmética de punteros y el bucle en K con
\texttt{tl.dot} y acumulación en fp32. Son propios el espacio de autotuning y la
comprobación en el PTX de que se usan tensor cores. \texttt{matmul\_tma.py} toma del
tutorial~09 los descriptores de tensor creados dentro del kernel
(\texttt{tl.make\_tensor\_descriptor}) y el \emph{allocator} obligatorio
(\texttt{triton.set\_allocator}). A diferencia del tutorial, no es persistente: conserva
la rejilla del baseline para aislar el efecto de TMA. \texttt{matmul\_blockptr.py} se
mantiene como intento documentado de activar TMA con \emph{block pointers}, que en
sm\_121 bajan a \texttt{cp.async} clásico (sección~\ref{sec:tma}). El matmul de Helion
es el del ejemplo oficial sin el \emph{epilogue}. El de CUTLASS ejecuta sin modificar el
ejemplo de CuTe DSL para Blackwell GeForce (\texttt{Sm120GemmKernel}: fp16, TMA y pipeline de
varias etapas). Solo es propia la envoltura que le pasa las matrices de PyTorch sin copiarlas.
No admite bf16.

\subsection{Kernels memory-bound: RMSNorm}

\begin{table}[htbp]
  \centering\small
  \begin{tabularx}{\linewidth}{lllX}
    \toprule
    kernel (\texttt{src/MBKernels/}) & DSL & elaboración & fuente \\
    \midrule
    \texttt{triton/rmsnorm\_baseline.py} & Triton & adaptado & tutorial de Triton \texttt{05-layer-norm}~[1] \\
    \texttt{triton\_tlx/rmsnorm.py} & Triton + TLX & propio & API de µTLX (\texttt{triton-ext})~[4] \\
    \texttt{gluon/rmsnorm.py} & Gluon & propio & API de Gluon, tutorial \texttt{gluon/02-layouts}~[1] \\
    \texttt{helion/rmsnorm.py} & Helion & adaptado & ejemplo oficial \texttt{examples/rms\_norm.py}~[3] \\
    \texttt{cutlass/rmsnorm\_ext.cu} & CUTLASS & oficial sin modificar & \texttt{cutlass::rmsnorm} (\texttt{device\_rmsnorm.h})~[2] \\
    \bottomrule
  \end{tabularx}
  \caption{Procedencia de las implementaciones de RMSNorm.}
  \label{tab:procedencia-mb}
\end{table}

El RMSNorm en Triton adapta el LayerNorm del tutorial~05: un programa por fila,
\texttt{MAX\_FUSED\_SIZE}~$=65536$ y la heurística de warps \texttt{BLOCK\_SIZE // 256}
(hasta 16 warps en lugar de 8). Se elimina la resta de la media. La versión Gluon repite
el algoritmo, pero fija el reparto de la fila entre hilos con un \texttt{BlockedLayout}
explícito. La de Helion es el \texttt{rms\_norm\_fwd} oficial sin la salida
\texttt{inv\_rms}, que solo necesita el \emph{backward}. La de CUTLASS invoca la función
oficial \texttt{cutlass::rmsnorm} sin modificarla: en fp16 usa el kernel vectorizado con
\texttt{float4} y en bf16 el escalar. Solo la envoltura como extensión de PyTorch es
propia. La versión TLX es un diseño propio (kernel persistente con \emph{prefetch} de
filas a memoria compartida mediante \texttt{cp.async}) que sigue el patrón de los tests y
tutoriales de triton-ext. Con \texttt{triton-utlx}~3.8.0.post1 no compila por un error del
plugin: crea \texttt{ttg.async\_copy\_global\_to\_local} sin \texttt{operandSegmentSizes},
y el fallo se reproduce también con el ejemplo mínimo del propio repositorio.

Como referencia de rendimiento y de corrección se usan implementaciones de PyTorch, no
kernels propios: \texttt{torch.matmul} (cuBLAS), \texttt{torch.\_scaled\_mm} (cuBLASLt,
FP8), \texttt{to\_sparse\_semi\_structured} (cuSPARSELt 2:4, con poda 2:4 propia) y
\texttt{torch.nn.functional.rms\_norm}.

\begin{table}[htbp]
  \centering\small
  \begin{tabularx}{\linewidth}{lXll}
    \toprule
    & proyecto & versión & licencia \\
    \midrule
    {[1]} & Triton (\texttt{triton-lang/triton}, \texttt{python/tutorials/}) & v3.8.0 (\texttt{c01b677}) & MIT \\
    {[2]} & CUTLASS (\texttt{NVIDIA/cutlass}) & v4.8.0 (\texttt{098de2a}) & BSD-3-Clause \\
    {[3]} & Helion (\texttt{pytorch/helion}) & \texttt{main} (\texttt{0f3b242}) & BSD-3-Clause \\
    {[4]} & µTLX (\texttt{triton-lang/triton-ext}, \texttt{triton-utlx} 3.8.0.post1) & \texttt{main} (\texttt{3c2bb3c}) & MIT \\
    \bottomrule
  \end{tabularx}
  \caption{Fuentes de los kernels. Todas las licencias permiten reutilizar y modificar el código conservando el aviso de copyright.}
  \label{tab:fuentes}
\end{table}
.
% Requiere: \usepackage{booktabs,tabularx}. Version Markdown: docs/TFM/procedencia_kernels.md.

\section{Procedencia de los kernels}
\label{sec:procedencia}

Esta sección indica el origen de cada kernel: si es código oficial sin modificar, una
adaptación de un ejemplo público o una implementación propia sobre una API documentada.
La procedencia se ha verificado comparando el código con la fuente citada, en la versión
indicada en la tabla~\ref{tab:fuentes}. Se distinguen tres tipos de elaboración:
\textbf{oficial sin modificar} (se llama al algoritmo de la fuente tal cual, con la
envoltura mínima para medirlo), \textbf{adaptado} (parte de un ejemplo público, con los
cambios que se detallan) y \textbf{propio} (diseño propio sobre la API documentada).

\subsection{Kernels compute-bound}

\begin{table}[htbp]
  \centering\small
  \begin{tabularx}{\linewidth}{llX}
    \toprule
    kernel (\texttt{src/CBKernels/triton/}) & elaboración & fuente \\
    \midrule
    \texttt{matmul.py} (baseline) & adaptado & tutorial de Triton \texttt{03-matrix-multiplication}~[1] \\
    \texttt{matmul\_blockptr.py} & adaptado & baseline + \texttt{tl.make\_block\_ptr}~[1] \\
    \texttt{matmul\_tma.py} & adaptado & tutorial de Triton \texttt{09-persistent-matmul} (variantes TMA)~[1] \\
    \texttt{matmul\_fp8.py} & adaptado & baseline con entradas FP8 e4m3~[1] \\
    \texttt{matmul\_fp8\_tma.py} & propio & combinación de \texttt{matmul\_fp8} y \texttt{matmul\_tma} \\
    \midrule
    \texttt{helion/matmul.py}$^\dagger$ & adaptado & ejemplo oficial de Helion \texttt{examples/matmul.py}~[3] \\
    \texttt{cutlass/matmul\_cute.py}$^\dagger$ & oficial sin modificar & CuTe DSL \texttt{blackwell\_geforce/kernel/dense\_gemm/dense\_gemm.py}~[2] \\
    \bottomrule
  \end{tabularx}
  \caption{Procedencia de los kernels de matmul. $^\dagger$Ruta relativa a \texttt{src/CBKernels/}.}
  \label{tab:procedencia-cb}
\end{table}

El baseline \texttt{matmul.py} comparte 27 de sus 36 líneas distintas con el
\texttt{matmul\_kernel} del tutorial~03. Coinciden el orden de programas agrupado para la
L2 (\texttt{GROUP\_SIZE\_M}), la aritmética de punteros y el bucle en K con
\texttt{tl.dot} y acumulación en fp32. Son propios el espacio de autotuning y la
comprobación en el PTX de que se usan tensor cores. \texttt{matmul\_tma.py} toma del
tutorial~09 los descriptores de tensor creados dentro del kernel
(\texttt{tl.make\_tensor\_descriptor}) y el \emph{allocator} obligatorio
(\texttt{triton.set\_allocator}). A diferencia del tutorial, no es persistente: conserva
la rejilla del baseline para aislar el efecto de TMA. \texttt{matmul\_blockptr.py} se
mantiene como intento documentado de activar TMA con \emph{block pointers}, que en
sm\_121 bajan a \texttt{cp.async} clásico (sección~\ref{sec:tma}). El matmul de Helion
es el del ejemplo oficial sin el \emph{epilogue}. El de CUTLASS ejecuta sin modificar el
ejemplo de CuTe DSL para Blackwell GeForce (\texttt{Sm120GemmKernel}: fp16, TMA y pipeline de
varias etapas). Solo es propia la envoltura que le pasa las matrices de PyTorch sin copiarlas.
No admite bf16.

\subsection{Kernels memory-bound: RMSNorm}

\begin{table}[htbp]
  \centering\small
  \begin{tabularx}{\linewidth}{lllX}
    \toprule
    kernel (\texttt{src/MBKernels/}) & DSL & elaboración & fuente \\
    \midrule
    \texttt{triton/rmsnorm\_baseline.py} & Triton & adaptado & tutorial de Triton \texttt{05-layer-norm}~[1] \\
    \texttt{triton\_tlx/rmsnorm.py} & Triton + TLX & propio & API de µTLX (\texttt{triton-ext})~[4] \\
    \texttt{gluon/rmsnorm.py} & Gluon & propio & API de Gluon, tutorial \texttt{gluon/02-layouts}~[1] \\
    \texttt{helion/rmsnorm.py} & Helion & adaptado & ejemplo oficial \texttt{examples/rms\_norm.py}~[3] \\
    \texttt{cutlass/rmsnorm\_ext.cu} & CUTLASS & oficial sin modificar & \texttt{cutlass::rmsnorm} (\texttt{device\_rmsnorm.h})~[2] \\
    \bottomrule
  \end{tabularx}
  \caption{Procedencia de las implementaciones de RMSNorm.}
  \label{tab:procedencia-mb}
\end{table}

El RMSNorm en Triton adapta el LayerNorm del tutorial~05: un programa por fila,
\texttt{MAX\_FUSED\_SIZE}~$=65536$ y la heurística de warps \texttt{BLOCK\_SIZE // 256}
(hasta 16 warps en lugar de 8). Se elimina la resta de la media. La versión Gluon repite
el algoritmo, pero fija el reparto de la fila entre hilos con un \texttt{BlockedLayout}
explícito. La de Helion es el \texttt{rms\_norm\_fwd} oficial sin la salida
\texttt{inv\_rms}, que solo necesita el \emph{backward}. La de CUTLASS invoca la función
oficial \texttt{cutlass::rmsnorm} sin modificarla: en fp16 usa el kernel vectorizado con
\texttt{float4} y en bf16 el escalar. Solo la envoltura como extensión de PyTorch es
propia. La versión TLX es un diseño propio (kernel persistente con \emph{prefetch} de
filas a memoria compartida mediante \texttt{cp.async}) que sigue el patrón de los tests y
tutoriales de triton-ext. Con \texttt{triton-utlx}~3.8.0.post1 no compila por un error del
plugin: crea \texttt{ttg.async\_copy\_global\_to\_local} sin \texttt{operandSegmentSizes},
y el fallo se reproduce también con el ejemplo mínimo del propio repositorio.

Como referencia de rendimiento y de corrección se usan implementaciones de PyTorch, no
kernels propios: \texttt{torch.matmul} (cuBLAS), \texttt{torch.\_scaled\_mm} (cuBLASLt,
FP8), \texttt{to\_sparse\_semi\_structured} (cuSPARSELt 2:4, con poda 2:4 propia) y
\texttt{torch.nn.functional.rms\_norm}.

\begin{table}[htbp]
  \centering\small
  \begin{tabularx}{\linewidth}{lXll}
    \toprule
    & proyecto & versión & licencia \\
    \midrule
    {[1]} & Triton (\texttt{triton-lang/triton}, \texttt{python/tutorials/}) & v3.8.0 (\texttt{c01b677}) & MIT \\
    {[2]} & CUTLASS (\texttt{NVIDIA/cutlass}) & v4.8.0 (\texttt{098de2a}) & BSD-3-Clause \\
    {[3]} & Helion (\texttt{pytorch/helion}) & \texttt{main} (\texttt{0f3b242}) & BSD-3-Clause \\
    {[4]} & µTLX (\texttt{triton-lang/triton-ext}, \texttt{triton-utlx} 3.8.0.post1) & \texttt{main} (\texttt{3c2bb3c}) & MIT \\
    \bottomrule
  \end{tabularx}
  \caption{Fuentes de los kernels. Todas las licencias permiten reutilizar y modificar el código conservando el aviso de copyright.}
  \label{tab:fuentes}
\end{table}
